% Generated draft: requires mathematical and source review.
\documentclass[11pt]{article}
\usepackage[margin=25mm]{geometry}
\usepackage{fontspec}
\setmainfont{Latin Modern Roman}
\usepackage{amsmath,amssymb,mathtools,unicode-math}
\setmathfont{Latin Modern Math}
\usepackage{enumitem,needspace,xurl,hyperref}
\setlength{\parindent}{0pt}
\setlength{\parskip}{5pt}
\setlength{\emergencystretch}{3em}
\newcommand{\N}{\mathbb{N}}
\newcommand{\Z}{\mathbb{Z}}
\newcommand{\Q}{\mathbb{Q}}
\newcommand{\R}{\mathbb{R}}
\newcommand{\C}{\mathbb{C}}
\newcommand{\F}{\mathbb{F}}
\newcommand{\A}{\mathbb{A}}
\newcommand{\PP}{\mathbb{P}}
\newcommand{\E}{\mathbb{E}}
\newcommand{\eps}{\varepsilon}
\newcommand{\dd}{\,\mathrm{d}}
\begin{document}
\section*{UnsolvedMath ID: 30000884}
\textbf{Source problem code:} catalogue entry
\subsection{Definitions and mathematical statement}
For a finite group $H$, let $\varepsilon:\Z H\to\Z$ be augmentation, $\varepsilon(\sum_h a_hh)=\sum_h a_h$, and put
\[
V(\Z H)=\{u\in(\Z H)^\times:\varepsilon(u)=1\}.
\]
A torsion unit is a unit $u$ such that $u^m=1$ for some integer $m\geq1$.

The first question has the quantifiers
\[
\forall H\text{ finite}\quad\exists G\text{ finite with }H\leq G\quad
\forall u\in V(\Z H)\text{ torsion}\quad
\exists x\in(\Q G)^\times\ \exists h\in H:\ x^{-1}ux=h.
\]
Here $\Z H\subseteq\Q G$ is induced by the chosen inclusion. A single overgroup $G$ must work for every torsion unit of $\Z H$. The source then asks the stronger rational-conjugacy question: can $x$ always be chosen in $(\Q H)^\times$? Such a choice would not depend on enlarging the ambient group algebra. Containment is written $H\leq G$ without imposing an unmentioned requirement that it be proper.

The second collection of questions concerns a finite group $G$, a prime $p$, a simple $\Q G$-module $X$, and a $\Z_{(p)}G$-lattice $L\subset X$, where
\[
\Z_{(p)}=\{a/b\in\Q:p\nmid b\}.
\]
Under the usual full-lattice convention, $L$ is finitely generated free over $\Z_{(p)}$, is $G$-stable, and satisfies $\Q\otimes_{\Z_{(p)}}L=X$. Determine the structure of
\[
E=\operatorname{End}_{\Z_{(p)}G}(L)
=\{a\in D:aL\subseteq L\},\qquad D=\operatorname{End}_{\Q G}(X),
\]
an order over $\Z_{(p)}$ in the rational division algebra $D$. The problem further requests the local-order building blocks arising from central primitive idempotents of $\Q G$, and a characterization of which such local orders can occur.

Because $X$ is simple, exactly one central primitive idempotent of $\Q G$ acts as the identity on $X$, and the others act as zero. Thus the phrase about building blocks is not by itself a nontrivial product decomposition of this particular $E$; the intended order construction, decomposition and equivalence notion are not specified. Likewise ``local order'' may refer to localization at $p$ rather than asserting that $E$ is a local ring. These meanings must be clarified for a complete classification statement.

The supplied triage checked 2026-08-21 labels the problem partially solved. It reports a positive answer for the prime-order restriction of the overgroup question and many rational-conjugacy cases, while calling the unrestricted finite-overgroup and endomorphism-ring questions unresolved. These are dated attributed claims, not conclusions proved here.
\subsection{Short English statement}
Ask whether one finite overgroup suffices to conjugate every normalized torsion unit of $\Z H$ to an element of $H$, and whether conjugation inside $\Q H$ suffices. Separately, determine the endomorphism orders of localized group lattices and classify their proposed local building blocks.
\subsection{Sources}
\begin{itemize}
\item Local UnsolvedMath registry, permanent ID 30000884, source code catalogue entry; supplied complete record. Its corrected-verified note says that all conjugacy and endomorphism-order questions were preserved and notation restored.
\item Mini-Workshop: Arithmetik von Gruppenringen (2008), Oberwolfach Reports, DOI 10.4171/owr/2007/55, pp. 3209--3240: \url{https://doi.org/10.4171/owr/2007/55}. The proposed-year field gives 2007.
\item Supplied rational-conjugacy reference: \url{https://arxiv.org/abs/1305.7419}.
\end{itemize}
\end{document}
